%% July 31, 2003
\documentclass[12pt]{article}
\usepackage{latexsym,graphicx,afterpage,epic,eepic,pstcol,comment}
% \setlength{\topmargin}{0.5in}
\setlength{\oddsidemargin}{0.25in}
\setlength{\textwidth}{6in}
\setlength{\textheight}{8.5in}

\newcommand{\Class}[1]{\textbf{#1}}
\newcommand{\Code}[1]{\texttt{#1}}
\newcommand{\Env}[1]{\texttt{#1}} %% LaTeX environment or package
\newcommand{\Ext}[1]{\texttt{#1}} %% file extension/type
\newcommand{\Program}[1]{\texttt{#1}}

\newcommand{\Compound}{\Class{Compound}}
\newcommand{\Compounds}{\Class{Compound}s}
\newcommand{\Camera}{\Class{Camera}}
\newcommand{\Label}{\Class{Label}}
\newcommand{\Labels}{\Class{Label}s}
\newcommand{\Mesh}{\Class{Mesh}}
\newcommand{\Meshes}{\Class{Mesh}es}
\newcommand{\Object}{\Class{Object}}
\newcommand{\Objects}{\Class{Object}s}
\newcommand{\Path}{\Class{Path}}
\newcommand{\Paths}{\Class{Path}s}
\newcommand{\Picture}{\Class{Picture}}
\newcommand{\Pictures}{\Class{Picture}s}
\newcommand{\Plot}{\Class{Plot}}
\newcommand{\Plots}{\Class{Plot}s}
\newcommand{\Point}{\Class{Point}}
\newcommand{\Points}{\Class{Point}s}
\newcommand{\Vector}{\Class{Vector}}
\newcommand{\Vectors}{\Class{Vector}s}

\newcommand{\C}{\Code{C}}
\newcommand{\CXX}{\Code{C++}}

\newcommand{\eepic}{\Env{eepic}}
\newcommand{\picenv}{\Env{picture}}

\newcommand{\eps}{\Ext{eps}}
\newcommand{\pdf}{\Ext{pdf}}

\newcommand{\elaps}{\Program{elaps}}
\newcommand{\ePiX}{\Program{ePiX}}
\newcommand{\epix}{\Program{epix}}
\newcommand{\pyepix}{\Program{Pyepix}}

%%\renewcommand{\LaTeX}{{\rm LaTeX}} %% For latex2html

\newcommand{\bi}{\begin{itemize}}
\newcommand{\ei}{\end{itemize}}

\newcommand{\bd}{\begin{description}}
\newcommand{\ed}{\end{description}}

\title{\ePiX~0.9.x Draft Specification}
\author{Andrew D. Hwang}
\date{Version 0.7, July 31, 2003}

\begin{document}

\maketitle

\section{Introduction}

\ePiX\ is a picture description library written in \CXX\ and designed
for creation of mathematically accurate line figures in \LaTeX. From a
user's perspective, \ePiX\ provides a command-driven front end to the
\LaTeX\ picture environment. Figure elements are described in
mathematically natural coordinates instead of picture coordinates, and
data from numerical computations is incorporated into the output. By
separating the abstract size of a figure from its printed size, \ePiX\ 
facilitates rescaling without rewriting the entire figure;
effectively, \ePiX\ is a vector format. Placement of text is
accomplished in a manner that is attractive over a large range of
printed scales.

The core library provides a simple, uniform, recursively structured
interface to basic geometric objects and plotting capabilities, while
providing hooks for user extensions and customization. The final
appearance of a figure is wholly user-specified; there are almost no
default behaviors. The guiding principles are logical structuring and
mathematical accuracy, typographical quality of output, ease of use,
and efficiency (run time and output file size).

The current design is heavily influenced by discussions with Robin
Blume-Kohout; in particular, the recursive structure of the \Picture\ 
class (which is both simplifying and enhancing) is Robin's idea, as is
the suggestion to move from a procedurally-oriented language to an
object-oriented one. Of course, all errors of design are entirely my
fault.


\subsection{Overview}

Figure elements live in a 3-D Cartesian \emph{object space}, and are
rendered in a 2-D \emph{screen plane}, which is mapped affinely onto
the printed page. Object space is populated by three types of figure
element: text labels, paths, and surface meshes. The screen plane
contains projected images of these elements (Figure~\ref{fig:camera})
and provides page mark-up by allowing multiple subfigures to be
amalgamated precisely (Figure~\ref{fig:import}).

The fundamental class is the \Picture, an abstraction that provides a
3-D coordinate system, camera (\emph{q.v.}), screen plane, and
numerous options that govern the style and appearance of the printed
figure, and that allow figure elements to be removed or truncated
according to precise, flexible criteria. The camera works as in
version~0.8.10 (see below), and determines the mapping from object
space to the screen plane.  Applications of the coordinate system
include logarithmic plotting, spherical and cylindrical coordinates,
and reflection of a figure through a plane.

An \ePiX\ figure comprises a tree of \Pictures\ with the entire figure
as root. sub\Pictures\ are designed for creation of compound figures,
by including several graphs in a single \LaTeX\ picture environment.
Screen coordinates of a \Picture\ are used to ``import''
the picture into its parent:

\begin{figure}[hbt]
\begin{center}
\input{import.eepic}
\caption{The unique affine map that imports a sub\Picture.}
\label{fig:import}
\end{center}
\end{figure}

Visible elements of a \Picture\ are \Objects, conceptually one of the
three types (label, path, or mesh) mentioned earlier. \Objects\ may be
grouped into \Compounds, or compound \Objects, and manipulated as a
single entity. A \Compound\ may contain other \Compounds, so within
a single \Picture, \Objects\ are organized as a collection of trees.
Unlike \Pictures, for which recursion depth greater than one seems not
to be useful, \Compounds\ of large depth have obvious applications.
Of course, it is an error for a \Compound\ to contain itself as an
\Object.

When a sub\Picture\ is imported, its \Objects\ are projected to the
screen plane and incorporated into the parent.  (\Compounds\ are
expanded recursively into lists of \Objects.)  Importing of
sub\Pictures\ therefore involves a loss of information for 3-D
elements, while inclusion of \Compounds\ does not.  Everything sent
to the output stream is an \Object.  When the root \Picture\ is
``imported'' (to the physical page), its \Objects\ are written to the
output stream.

An \ePiX\ input file consists of \CXX\ definitions, attribute-setting
and changing commands, and figure elements. Definitions include
user-defined constants, numerical functions, algorithms, and built-in
geometric data structures. Attribute commands change default behavior,
and remain in effect until superceded (or until the \Picture\ ends).


\section{The \Picture\ class}

The \Picture\ class is the primary means of organization within a
figure. As such, it is by far the most complicated \ePiX\ class.

\subsection{Coordinate systems}

Each \Picture\ has three coordinate systems associated to it, which
are described in turn below. These coordinate systems and the mappings
between them may be visualized as follows:

\begin{figure}[hbt]
\begin{center}
\input{coordinates.eepic}
\caption{\ePiX's coordinate systems}
\label{fig:coordinates}
\end{center}
\end{figure}


\medskip

\noindent\textbf{Abstract coordinates} are the 3-D coordinate system
in which \Picture\ elements are specified. Except to specify the
actual size of a \Picture, users interact only this coordinate system.


\medskip

\noindent\textbf{Object coordinates} are the 3-D Cartesian coordinate
system in which \Picture\ elements are represented and stored. Their
purpose is to provide a uniform coordinate system for all \Picture\ 
elements. Of course, it is necessary to store objects in spatial
coordinates until the \Picture\ is imported to the parent, to permit
hidden object removal and other mangling.

The \textbf{clip region} is a subset of object space, by default a
very large box, that is used to remove objects whose coordinates are
``too large'' (for example, too large for \LaTeX\ to handle).
\Objects\ are clipped as they are created rather than when the
\Picture\ is imported; unsetting the ``clip'' attribute does not
restore clipped portions of a figure. The \Picture\ class provides an
operator to clip all current elements.

The \textbf{User mapping} is, by default, the identity mapping, though
spherical, cylindrical, and various logarithmic and log-linear
coordinates are provided. Further, the mapping can be defined by the
user, for example to depict reflection through a plane, inversion in a
sphere, or rotation about an axis. Once a \Picture\ element has been
specified, a subsequent change of User mapping has no effect, since
elements are stored in Object coordinates.


\medskip

\noindent\textbf{Screen coordinates} are a 2-D (conceptually
Cartesian) coordinate system in which \Picture\ elements are rendered.

Analogous to the clip box is the \textbf{crop mask}, a region in the
screen plane outside (or inside) of which objects can be removed. The
crop mask is inscribed in a rectangle called the \emph{crop box}, and
may be rectangular (with rounded corners of specified radius),
diamond-shaped, or elliptical:

\begin{figure}[hbt]
\begin{center}
\input{crop_masks.eepic}
\caption{\ePiX's cropping masks}
\label{fig:crop_masks}
\end{center}
\end{figure}

The \textbf{Camera} (\emph{q.v.}) provides a mapping (the
\textbf{lens}) from object space to the screen plane. The camera is
meant to be flexibly controlled by the user, similarly to a real
camera. The default lens is orthogonal projection along the third
coordinate, which provides a sensible default for 2-D figures.


\medskip

\noindent\textbf{Ambient coordinates} do not properly belong to a
\Picture, but instead are screen coordinates of its parent. Ambient
coordinates of the root \Picture\ are \LaTeX\ picture coordinates,
which specify location on the printed page.

The \textbf{picture box} is a rectangle in ambient coordinates whose
location is where the parent ``expects'' the sub\Picture\ to be
placed. For the root \Picture, the picture box is the \LaTeX\ size and
offset of the figure in terms of the current \LaTeX\ 
\texttt{unitlength}, and corresponds to the (typographical) box which
\LaTeX\ has allotted to the figure. A \Picture\ does not ``know'' its
picture box.

A \Picture\ has a \textbf{bounding box}, a rectangle in the screen
plane that determines the default crop mask and size of the \Picture.
The \textbf{Import} mapping is the unique affine scaling between the
bounding box and the picture box. When a \Picture\ is imported, each
of its elements is rendered in 2-D, then ``pasted'' into the screen
plane of the parent, Figure~\ref{fig:import}.  Importing is intended
(only) to facilitate ``abstract mark-up'', namely to allow multiple
\Pictures\ to be incorporated as typographical subfigures.


\subsection{Style Attributes}

The final appearance of a figure (color, line widths, shape and style
of the border, etc.) can be customized easily. The only defaults are
those that seem sensible for simple, 2-D figures without
\texttt{PSTricks}: Black printing in a vanilla picture environment,
with the $(x,y)$-plane seen from straight down the $z$-axis.
Especially, \ePiX\ does not automatically draw bounding boxes or
coordinate axes, crop, clip, choose aspect ratios, or otherwise try to
make aesthetic choices regarding appearance.

\subsubsection*{Clipping and Cropping}

Clipping is a Boolean attribute in object space that affects whether
or not an \Object\ is rendered, and that determines what parts of a
\Path~or \Mesh\ are visible. As mentioned above, each \Picture\ has a
clip region, implemented as a Boolean function of position called
\verb+is_visible+. By default, the clip region is a very large box
centered at the origin; functionally, clipping is off.  (Literal
clipping always occurs, to eliminate points whose coordinates are too
large for \LaTeX\ to handle.)  Clipping is applied to \Objects\ when
they are created (clipped portions of \Objects\ are not stored), or by
explicit request, in which case a clipped \Object\ is removed from the
\Picture, and a \Path~or \Mesh\ is appropriately truncated.

Cropping is a Boolean attribute in the screen plane that masks
elements when a \Picture\ is imported. Cropping is off by default, and
is not meaningful for individual elements, but only for the \Picture\ 
as a whole. If it is desirable that only some elements be cropped, two
\Pictures\ may be superimposed and differently cropped.

Clipping and cropping may be done either including or excluding the
boundary of the respective region (mathematically, to clip or crop to
a closed or open set).  The bounding box or crop mask of a \Picture\ 
is an ordinary element of type \Path\ (\emph{q.v.}), but has zero
width (i.e., is not drawn) by default.

\subsubsection*{Color}

Colors in an \ePiX\ figure are given in either an RGB or a CMYK
specification, as in earlier versions. By default, the scheme is RGB.
The color scheme should not be changed within a \Picture; there is no
reason to, and \ePiX\ discourages it. Colors are implemented via
\texttt{PSTricks}, which involves \texttt{special}s in the DVI file.

A \Picture\ may be either opaque or transparent; in implementation,
the difference is whether or not a solid rectangle is placed under the
\Picture's contents. The background color is meaningful only when a
figure is opaque, and is enacted when the \Picture\ is exported.

The foreground color of a \Picture---the color of its
elements---defaults to black for the root \Picture, and is inherited
by sub\Pictures. Elements may specify an explicit color, or may be
colored according to the default.


\subsection{Class description}

\noindent The \ePiX\ \Picture\ class has the following members and
functions:

\bi
  \item Camera (\emph{q.v.})

  \item User coordinate system
    \bi
      \item \verb+set_user_coords(triple user_coords(triple))+
    \ei

  \item \verb+bounding_box+, canonicalized screen rectangle
    \bi
      \item \verb+bounding_box(triple p1, triple p2)+. Only the first
      two components of~$p_i$ are meaningful.
    \ei

  \item Clip region: \verb+bool is_visible(triple)+
    \bi
      \item \verb+clipping(bool)+. Turn clipping on/off.
        
      \item \verb+clip_to(bool is_visible(triple))+ is the most
        general way of setting clip region, and allows an arbitrary
        Boolean function to be specified.  In all cases, a
        user-specified clip mask is ``\texttt{and}''ed with the
        default mask, which checks whether all coordinates are smaller
        than \verb+EPIX_INFTY+ in absolute value.

      \item \verb+clip_box(triple p1, triple p2)+. Use specified
        arguments as corners of the clip box.

      \item \verb+clip_box(triple p1)+. Use $\pm$~specified
        argument as corners of the clip box.

      \item \verb+clip_to(triple p1)+. Use origin and specified
        argument as corners of the clip box.
        
      \item \verb+clip_to(triple p1, double radius)+. Set clip region
        to be a ball with given center and radius.

      \item \verb+set_clipping(OPEN/CLOSED)+. Use strict or non-strict
        inequality when testing against the clip box or ball.

      \item \verb+clip_all()+. Clip all current \Picture\ elements. 

      \item Possible implementation: Maintain two global data
        structures, a box and ball, which by default are centered at
        the origin and have diameter \verb+2*EPIX_INFTY+. Define
        Boolean functions that read these data structures, say
        \verb+clip_mask_box+~and \verb+clip_mask_ball+, and have each
        user interface function set the \verb+clip_mask+ to be an
        appropriate function.
    \ei

  \item Crop mask (defaults to bounding box)
    \bi
      \item \verb+crop_box(triple p1, triple p2)+. Set circumscribing
        rectangle with specified corners, relative to \texttt{target}
        (only first two components significant).

      \item \verb+crop_box(triple p1)+. Set circumscribing rectangle
        with $\pm$~specified corner, centered on \texttt{target} (only
        first two components significant).

      \item \verb+crop_to(double radius)+. Set circumscribing
        rectangle as above, set crop mask to \Code{ELLIPSE}.

      \item \verb+crop_mask(EPIX_CROP_MASK)+. (\Code{BOX},
        \Code{DIAMOND}, or \Code{ELLIPSE}).
      
      \item \verb+set_cropping(OPEN/CLOSED)+. Use strict or non-strict
        inequality when cropping.

      \item \verb+border_width(double)+,  \verb+border_width(dimen)+.
        Set the width of the border (same as \Code{pen} function).

      \item \verb+border_style(EPIX_PATH_STYLE)+ (solid, dashed, or
        dotted).

      \item \verb+border_color(color spec)+.

    \ei

  \item Color and style
    \bi
      \item \Code{pen()}, \Code{plain()}, \Code{bold()}. Same as in
        old version.

      \item \verb+path_style(EPIX_PATH_STYLE)+, \Code{solid()},
        \Code{dashed()}, \Code{dotted()}. Same as in old version.

      \item \verb+color_scheme(RGB/CMYK)+. Defaults to \Code{RGB}.

      \item \Code{decolorize()}. Strip color statements from output.

      \item \verb+background_color(color spec)+,
        \verb+foreground_color(color spec)+. Default is parent's
        value.

      \item \verb+transparent(bool)+. (Default is \Code{true}.)

      \item \verb+fill_style(fill spec)+. Set fill style for closed
        \Paths.

      \item \verb+fill_color(color spec)+. Set fill color for closed
        \Paths.
    \ei

  \item Text
    \bi
      \item \verb+text_angle(angle)+. Rotates text labels by specified
        angle. Same effect can be achieved manually by including the
        rotate statement into label text. Effective only with
        \Code{pstcol}. 

      \item \verb+text_font(font spec)+. Set the font; default is
        NULL. Need not be implemented, since effect is easily achieved
        manually.
    \ei

  \item Linked list of \Picture\ elements
    \bi
      \item \Code{marker()}, \Code{label()}, \Code{plot()}. Add
        \Objects.

      \item \verb+clip_all()+. Clip current list using current clip
        region. 

      \item \verb+remove_hidden(bool)+. Set hidden \Object\ removal
        attribute (default is \Code{true}). Enacted at import.
        
      \item \Code{import()}, \Code{write()}. \Code{import} renders all
        elements in the screen plane, then maps them to the parent by
        affine scaling (Figure~\ref{fig:import}). sub\Pictures\ 
        are handled recursively. \Code{write} is similar, but writes
        to the output stream.

    \ei
\ei

\noindent Remarks:

\bi
  \item The camera defaults to current camera, user coordinates
    default to identity map, clip box defaults to ``large'' box, crop
    mask defaults to bounding box, clipping and cropping are off,
    transparency is on, background color is white, foreground color is
    black, border is empty/off.
    
  \item Cropping of top-level figure elements is done at the export
    stage, to ensure that a figure exactly fills the region allotted
    by the parent. Other globally-applied attributes (such as removal
    of all color statements) are also effected at this stage.
    
    \Picture\ elements may be automatically ordered in the output file
    according to various criteria, such as by color or line width.
    Explicit writing can be used to force a particular ordering in the
    output. 
\ei


\section{Other Classes}

The current design calls for almost 2~dozen auxiliary classes, ranging
from typographical attributes (e.g., color) to abstract mathematical
data types (lines, circles, spheres, etc.). The primary goal is to
build a modular, flexible library from simple components. The
dependency tree is shown in Figure~\ref{fig:classes}.
%%\afterpage{clearpage}
\begin{figure}[!hbt]
\begin{center}
\input{classes.eepic}
\caption{\ePiX's auxiliary classes}
\label{fig:classes}
\end{center}
\end{figure}

\subsection{Abstract Classes}

The first group of classes comprises \ePiX's internal types, used to
specify abstract data structures.

\subsubsection*{Position, Vector, and Frame}

The \Class{Position} class is a triple of numbers, namely a point of
object space. Mathematically, a Position is an element of Cartesian
3-space. A Position represents nothing but position; it does not make
sense to add Positions, or to multiply a Position by a scalar. In
other words, Positions should be viewed as coordinates on a manifold,
but not generally as affine/linear coordinates on Cartesian space; it
is only permissible to perform affine operations when the User
coordinate system is the identity map.

A \Class{Vector} is an infinitesimal displacement between Positions,
as in differential geometry, and is implemented as a Position (the
\emph{base}) together with an ordered triple of numbers (the
\emph{displacement}). Vector space operations are defined (only) for
Vectors based at the same Position. The following are defined:
(In)equality; addition and scalar multiplication; dot, cross, and
componentwise product; orthogonalization (Gram-Schmidt). The
constructors for Positions and Vectors are sensitive to the current
User coordinate system.  Non-Cartesian figures can be drawn by
defining suitable analogues of the dot product and orthogonalization.

When the User coordinate system is Cartesian, a vector may be viewed
as a literal displacement, and vectors based at different points may
be used as operands by parallel translating them to the same
basepoint.

A \Class{Frame} is a right-handed orthonormal triple of Vectors based
at a single point. A Frame may be rotated about any of its elements.
The primary use of a Frame is as the orientation of the Camera.

\subsubsection*{Pair}

A \Class{Pair} represents an element of the screen plane. The standard
constructor is Cartesian, and no other constructors or operations are
provided. 

\subsubsection*{Color\_spec}

A \Class{Color\_spec} is a set of color densities, both RGB~and CMYK.
Each figure has a global color model, and only the corresponding
densities are meaningful.

\subsubsection*{Camera}

The \Camera\ is implemented as in \ePiX-0.8.10: There is an
orthonormal frame labelled \textbf{sea}, \textbf{sky}, and
\textbf{eye} that consists of a right-pointing vector, an upward
pointing vertical vector, and their cross product. The Camera is
oriented by rotating about these axes.  The Camera ``looks at'' a
point in object space called the \textbf{target}, whose distance from
the Camera is the \textbf{range}. The plane through the target and
parallel to the sea-sky plane is the screen plane described above.
Finally, the mapping from object space to the screen plane is the
\textbf{lens}.  Lenses include point projection (simple perspective)
and fisheye (radial projection to a hemisphere, followed by axial
projection to the equatorial disk). Point projection from a point at
infinity is orthogonal projection.%  \medskip

\begin{figure}[hbt]
\begin{center}
\input{camera.eepic}
\caption{\ePiX's camera model}
\label{fig:camera}
\end{center}
\end{figure}

The default is for the Camera to look down the $z$-axis at the
$(x,y)$-plane with the axes in the standard orientation, and with the
lens being orthogonal projection. 

All \textbf{lens}es perform Euclidean point projection centered at the
Camera, then (possibly) do further transforming to get a planar image.
Lenses that do not perform Euclidean point projection will not work
with the default visibility algorithms, nor with the \textbf{draw}
functions of geometric types. However, internal algorithms that make
assumptions about the \textbf{lens} should provide hooks for users who
add ``non-Euclidean'' \textbf{lens}es.

\subsubsection*{Shape}

\ePiX\ provides a virtual \Class{Shape} class, and the following
built-in geometric types: line and segment, plane, circle, sphere.
When appropriate, geometric types can be scaled or translated. Each
type has a \textbf{draw} function that creates a \Path\ (\emph{q.v.}).
Geometric types may be intersected generically, in the sense that the
intersection of two ``randomly chosen'' objects is another object of
one of these classes, e.g., the intersection of a plane and a sphere
is a circle. Appropriate intersection operators are provided. The
following exceptions are thrown for intersections that are not
generic:

\bi
  \item Parallel, skew (for lines, planes), disjoint (for segments
  lying on lines that otherwise intersect)

  \item Non-coplanar, concentric, separated, tangent (circles, lines,
    and spheres)

  \item Coincindent (all types)
\ei

In addition to the ``obvious'' constructors, the following are
provided: Plane and circle constructors for three non-collinear
points, sphere constructor for four non-coplanar points, plane
constructor for a line and point.


\subsection{The \Object\ and \Compound\ classes}

The classes described above are internal implementation details; none
of them actually appears in a completed figure.  Visible elements are
\Objects, which have a location, color, contents (such as the text of
a \Label), and attributes (such as the rotation angle of a \Label, or
the width or style of a \Path). The \Objects\ of a \Picture\ are
stored as a doubly linked list to facilitate sorting and removal.

A \Compound\ is a doubly linked list of \Objects~and \Compounds. A
\Compound\ has an origin but knows nothing of its (abstract or
concrete) size. A \Compound\ may be translated, scaled, or rotated.
When placed in a \Picture, the \Compound's origin is translated to the
specified location. Finally, a \Compound\ possesses simple style
attributes: line width, path style, and color. These may be overridden
by individual elements, naturally.

\subsubsection*{Point (P)}

A \Point\ consists of a Position, Color\_spec, and a visibility flag.
By default, a~P is visible, but hidden-object removal algorithms work
(in part) by marking~Ps as invisible.

\subsubsection*{Label~and Marker}

A \Label\ is a \Picture\ element that, for \LaTeX\ purposes, is a
typographical box; this includes point markers (which are implemented
as boxes of size zero) and \LaTeX\ typography (``labels'') aligned on
a basepoint.  The text of a Label may be masked (with white by
default), which hides portions of the figure ``underneath'' the Label.
A Label consists of a Point, a screen plane offset (in true~pt), an
alignment option, and text content. 

A \Class{Marker} is a label of offset zero, aligned on its Point. The
``text'' of a Marker is an \verb+EPIX_MARK_TYPE+ enum.


\subsubsection*{Vertex}

The geometric workhorse for plots is the \Class{Vertex}, a~P together
with its distance to the Camera and a list of neighbors, each of which
abstractly represents a \Path\ segment.

\subsubsection*{Path}

A \Path\ is a list of Vertices representing a polygon, curve object, or
function plot. The \Path\ class has the following members:

\bi
  \item List of Vertices.

  \item Concrete style attributes: These attributes are effected by
    printing a command to the output file, and include line width,
    path style (solid, dashed, dotted), color, and whether or not a
    path is to be filled.
    
    For simulation of transparency (and other effects), it may be
    desirable for each \Path\ segment to have separate color, style,
    and width attributes.
\ei

\noindent The \Path\ class has the following functions:

\bi
  \item Constructors for ellipse, polygon/polyline, spline, and
    function plot

  \item Functions to set style parameters (width, style, color,
  closing, filling, fill style (hatching, color, etc.))

  \item Functions to split, reverse, and amalgamate \Paths
    
\ei

It may prove convenient to make \Path\ data (the list of Vertices) a
separate class. Note that a \Path\ does not know how to draw itself.

\subsubsection*{Face}

A \Class{Face} is an opaque triangle, usually one whose vertices are
adjacent Vertices of a Mesh (\emph{q.v.}). Faces are not usually drawn
(PostScript has a limited color stack, so using colored or shaded
Faces is potentially difficult), but are necessary for hidden line
removal. 

\subsubsection*{Mesh}

A \Mesh\ represents a piece of surface, and is stored as a double list
of \Paths~and a list of Faces. The \Paths' typographical attributes
are used to draw the actual grid lines, while the Faces are used
(only) for hidden line removal. Like a Path, a Mesh does not know how
to draw itself. Instead, a Picture being rendered performs hidden
object removal, then draws every Path and Mesh it contains.

The rough idea of hidden line removal is to compare every point
(Point, Label, Vertex) of a figure with every Face, to remove points
that are hidden, and to add in new points, as when part of a Path
segment is occluded and one endpoint must be moved. Brute force
entails an unacceptably long run time. The following outline is
something of an improvement.

When a Picture is exported, it contains a list of Labels, Paths, and
Meshes. Each Mesh is divided into \emph{layers}---regions where the
faces point toward the Camera or away from it. (The Mesh is oriented
by its parameter domain.) Each layer is mapped homeomorphically to the
screen, and the list of vertices in the screen plane is stored. Next,
each point is tested for occlusion by each layer; first a crude test
(is the image of the point inside the boundary of the screen image of
the layer?) is used to mark vertices as unhidden. When this test is
inconclusive (the image lies inside the boundary), the faces
comprising the layer are tested one by one, and occluding faces are
recorded. At the end of the process, each vertex in the Picture is
marked as hidden or not, and for each hidden vertex there is a list of
faces that occlude the vertex.

Figure elements that are point-like (Points and Labels) and hidden are
dropped at this stage. Next, Paths are culled: For each pair of
adjacent vertices with exactly one visible, a new visible vertex is
added as close to the hidden vertex as possible. (The list of faces
that occlude the hidden vertex is used to shorten the computation of
the new vertex.) When the new vertex is located, the hidden vertex
and all subsequent consecutive hidden vertices are removed, and the
Path is split into two Paths. At the end of the process, the figure
contains a list of Labels, Points, and Paths, none of which occludes
any other. 

Finally, the remaining elements are sorted by color (to minimize the
load on the color stack), then written to the output. The run time
should be considerably shorter than brute force, especially in usual
situations, for which each mesh has few layers.

\begin{comment}
Hidden line removal is effected when a \Picture\ is imported.  A
\Mesh\ is stored, manipulated, and printed as a linked list of \Paths,
but is treated as a collection of triangles for hidden line removal.
Each \Label~and \Path\ segment is tested for occlusion against every
\Mesh\ triangle in the \Picture. When a \Picture\ is imported, it
contains a linked list of elements: \Labels, \Paths, and \Meshes. For
each \Label\ location and \Path~or \Mesh\ vertex, test successively
for occlusion against each \Mesh\ triangle. If a \Label\ location is
hidden, the \Label\ is removed from the figure and testing begins with
the next element. For \Path\ segments, test for occlusion of one or
both endpoints. If both endpoints are hidden, the segment is removed
and testing resumes with the next segment. If only one endpoint is
hidden, a new vertex (on the boundary of the occluding triangle) is
computed, and the \Path\ is split into two \Paths. \Meshes\ are tested
as lists of \Paths. Note that \Meshes\ are treated differently as
elements to be rendered or as objects that hide other elements.

If a \Picture\ contains $k$~\Meshes\ of size~$n\times n$, then brute
force comparison gives an upper bound of $2k^2n^2(n+1)^2$~occlusion
tests. Even at 1~$\mu$sec per~test, a \Picture\ containing just
2~\Meshes\ each of size~$60\times60$ could take more than a minute to
render.  Brute force is unacceptably slow, though the number of tests
is likely to be smaller in practice. It is extremely desirable (and
worth a fairly large investment of algorithm design and code
optimization!) to find fast partial occlusion tests for entire \Path\ 
components. (For example, instead of testing a single segment against
all triangles, it may well be better to test an entire \Path\ 
component against a triangle, splitting into sub\Paths\ as required.)
\end{comment}

\end{document}


\subsection{Other functionality}

Compound objects, such as coordinate axes and grids, may be
implemented as collections of \Paths~and \Objects, though it may prove
fruitful to design special classes for these types, as well.

\section{Hiding}

For a 2-D figure done in black and white, layering and hidden object
removal are essentially non-existent issues. Colored planar figures
present a minor problem, though the solution is likely to be as simple
as ordering objects by color (darkest first?) as they are written to
the output file. 

The greatest unsolved problem in the current incarnation is hidden
object removal, which is prerequisite to full 3-D capability. One
simple but inelegant and inefficient algorithm (so far as eepic files
are concerned) is to build \Paths\ and surface meshes from line
segments and mesh quadrilaterals, to order these elements by the
distance from the camera to their barycenters, then to print the
elements in reverse order (farthest to nearest), letting Postscript
handle the resultant Z-buffering. There are at least two problems with
this idea: Quadrilaterals that intersect are not handled well, and a
potentially large number of hidden objects are written to the output
file. On the plus side, the results of fairly simple-minded \Object\ 
generation can look acceptably good:

\begin{figure}[hb]
\begin{center}
\includegraphics[width=4.5in]{juicer.eps}
\caption{A $48\times48$ surface mesh}
\label{fig:juicer}
\end{center}
\end{figure}

The file size of such an image is acceptable after conversion to
PDF, while the \eepic~and EPS files are large: In this example, the
\eepic\ file is 339,567~bytes, the EPS file is 449,729~bytes, and the
PDF file is 30,576~bytes.

-----
CLASS STRUCTURE

---
Abstract bounding_box class

Data:
* Cartesian coordinates of lower left and upper right corners
* Basepoint (default: lower left)
* Coordinate system (default: Use coordinate system of containing box)

Functions:
* Return data elements

---
Picture class

Data:
* bounding_box
* Relative size (wrt coordinates of containing Picture, if any)
* Absolute size (overrides relative size)
* Shape (rectangle, circle, oval, ellipse, diamond)
* Style parameters: Foreground color, border, masking, clipping
* List of Objects

Functions:
* Return size attributes
* Sort Objects by distance, color
* Cull hidden Objects, crop paths
* Draw 

---
Label class

Data:
* Basepoint
* Offset
* Orientation (rotation angle)
* Color, masking
* LaTeX command(s)

Functions:
* Return location (basepoint + offset in appropriate units)
* Draw

(A Marker is a Label with no offset or rotation, just a box of size
zero containing a LaTeX command.)

---
Path class

Data:
* Attributes (Color, path style and width, bool closed, shading for
  closed Paths) 
* Number of points, and list of points

Functions:
* Amalgamate (concatenate multiple Paths into one)
* Cull points (e.g., for object hiding), split Paths
* Draw (need to control file width, and number of points per \path to
  avoid LaTeX memory overflow)

(Polygons, arcs, ellipses, splines, and arrows will be implemented as
Paths.)

---
Face class

Data:
* Vertices

Functions:
* Orientation from camera ((counter)clockwise)
* Whether a point is hidden by the face
* Point(s) where segment hits hidden region


-----
HELPER CLASSES

Point class

Data:
* Cartesian coordinates

Functions: 
* Component functions, projections to coordinate axes/planes
* (In)equality operators
* Midpoint (convex linear combination)

---
Pair class (data structure for screen plane)

Data:
* Cartesian coordinates (type <complex>)

Functions:
* Dot product, multiplication by i

---
Vector class

Data:
* Basepoint (of type Point)
* Components

Functions:
* Component and basepoint functions
* Algebraic operations for vectors based at the same point (vector
  space operations; dot, cross, and componentwise products;
  (in)equality and orthogonalization operators)
* Polar constructor

---
Camera class

Data:
* Orientation (orthonormal basis {sea, sky, eye})
* Target (Point to look at)
* Distance to Target
* Lens (choice of projection mappig to the screen plane)

Functions:
* Rotate about sea, sky, or eye axes
* Get/change Target, Distance, or Lens
* Take a picture (project a point)